\documentclass[11pt, letter, english]{amsart}
\usepackage[margin=1in]{geometry}
\usepackage{amsmath} % utilities for mathematics
\usepackage{amsthm} % utilities for theorem environment
\usepackage{amssymb} % loads mathematical fonts
\usepackage[dvipsnames]{xcolor}
\usepackage{amscd} % utilities for commutative diagrams

\usepackage{mathtools}

\usepackage{subcaption}
\usepackage{array} % core package
\usepackage[sc]{mathpazo} % Palatino font with small caps
\usepackage[scr=boondoxo]{mathalfa}

\usepackage[unicode=true]{hyperref} % permits navigating on the pdf
\hypersetup{
    colorlinks=true,
    linkcolor=NavyBlue,
    citecolor=OliveGreen,
    urlcolor=blue,
    pdfencoding=auto
}
\usepackage{cleveref} % smart referencing
\usepackage{caption}

\usepackage{graphics,graphicx} % permits putting images
\usepackage{tikz,tikz-cd} % tool for diagrams

%\usetikzlibrary{graphdrawing.force}% quotes library is for the [""] edges
\usetikzlibrary{graphs,graphs.standard,calc}

\usetikzlibrary{decorations.pathreplacing,}%angles,quotes

\usepackage{enumerate} % permits enumerating using letters or other symbols

\DeclareMathAlphabet{\mathsf}{OT1}{\sfdefault}{m}{n}

\newcommand{\nocontentsline}[3]{}
\newcommand{\tocless}[2]{\bgroup\let\addcontentsline=\nocontentsline#1{#2}\egroup}

%\usepackage[margin=1.30in]{geometry}
%\usepackage[textwidth=360pt,textheight=615pt]{geometry}
\linespread{1.1}
\setlength\parindent{0.2in}

\usepackage{verbatim}
\usepackage{scalerel}

\makeatletter
\def\dual#1{\expandafter\dual@aux#1\@nil}
\def\dual@aux#1/#2\@nil{\begin{tabular}{@{}c@{}}#1\\#2\end{tabular}}
\makeatother

\makeatletter
\@namedef{subjclassname@2020}{\textup{2020} Mathematics Subject Classification}
\makeatother


\newcommand{\stirlingtwo}[2]{\biggl\{\genfrac{}{}{0pt}{}{#1}{#2}\biggr\}}
\newcommand{\tstirlingtwo}[2]{\left\{\genfrac{}{}{0pt}{}{#1}{#2}\right\}}
%\renewcommand{\tilde}{\widetilde}
%\renewcommand{\widetilde}{\wwtilde}

\DeclareMathAlphabet{\amathbb}{U}{bbold}{m}{n}


\newcommand{\mfrak}[1]{\scaleobj{0.93}{\mathfrak{#1}}}

\hypersetup{
    colorlinks = true,
    linkbordercolor = {white},
    linkcolor = {BrickRed},
    anchorcolor = {black},
    citecolor = {BrickRed},
    filecolor = {cyan},
    menucolor = {BrickRed},
    runcolor = {cyan},
    urlcolor = {black}
}

\usetikzlibrary{automata}
\usepackage{bbm}

\newtheoremstyle{teoremas}% <name>
{11pt}% <Space above>
{11pt}% <Space below>
{\itshape}% <Body font>
{}% <Indent amount>
{\bfseries}% <Theorem head font>
{}% <Punctuation after theorem head>
{.5em}% <Space after theorem headi>
{}% <Theorem head spec (can be left empty, meaning `normal')>



\theoremstyle{teoremas}
\newtheorem{theorem}{Theorem}[section]
\newtheorem{corollary}[theorem]{Corollary}
\newtheorem{lemma}[theorem]{Lemma}
\newtheorem{proposition}[theorem]{Proposition}

\newtheoremstyle{definition}% <name>
{11pt}% <Space above>
{11pt}% <Space below>
{}% <Body font>
{}% <Indent amount>
{\bfseries}% <Theorem head font>
{}% <Punctuation after theorem head>
{.5em}% <Space after theorem headi>
{}% <Theorem head spec (can be left empty, meaning `normal')>

\theoremstyle{definition}
\newtheorem{definition}[theorem]{Definition}
\newtheorem{thmdefinition}[theorem]{Theorem-Definition}

\newtheorem{conjecture}[theorem]{Conjecture}
\newtheorem{openproblem}[theorem]{Open Problem}
\newtheorem{problem}[theorem]{Problem}
\newtheorem{question}[theorem]{Question}
\newtheorem{example}[theorem]{Example}
\newtheorem{remark}[theorem]{Remark}

\crefname{theorem}{theorem}{theorems}
\Crefname{theorem}{Theorem}{Theorems}
\crefname{lemma}{lemma}{lemmas}
\Crefname{lemma}{Lemma}{Lemmas}
\crefname{proposition}{proposition}{propositions}
\Crefname{proposition}{Proposition}{Propositions}

\DeclareMathOperator{\rk}{rk}
\DeclareMathOperator{\cusp}{cusp}
\DeclareMathOperator{\Rel}{Rel}
\DeclareMathOperator{\Poin}{Poin}
\DeclareMathOperator{\Tor}{Tor}
\DeclareMathOperator{\Ext}{Ext}

\newcommand{\M}{\mathsf{M}}
\newcommand{\rev}{\operatorname{rev}}
%\newcommand{\B}{\mathsf{B}} let us only write \U_{n,n} for boolean matroids
\newcommand{\N}{\mathsf{N}}
\newcommand{\U}{\mathsf{U}}
\newcommand{\cI}{\mathcal{I}}
\newcommand{\A}{\mathrm{A}}
\newcommand{\Q}{\mathbb{Q}}
\newcommand{\R}{\mathbb{R}}
\newcommand{\Z}{\mathbb{Z}}
\newcommand{\LL}{\mathsf{\Lambda}}
\newcommand{\PP}{\mathsf{\Pi}}
\newcommand{\C}{\mathsf{C}}
%\renewcommand{\lambda}{\lambda}

\newcommand{\Hilb}{\operatorname{Hilb}}
\newcommand{\rank}{\operatorname{rk}}
\newcommand{\corank}{\operatorname{crk}}
\newcommand{\cl}{\operatorname{cl}}
\renewcommand{\H}{\mathrm{H}}
\newcommand{\CH}{\mathrm{CH}}
\newcommand{\aug}{\operatorname{aug}}
\newcommand{\IH}{\mathrm{IH}}
\newcommand{\uH}{\underline{\mathrm{H}}}
\newcommand{\uCH}{\underline{\mathrm{CH}}}
\newcommand\size[1]{|#1|}
\DeclareMathOperator{\Pal}{\mathrm{Pal}}
\DeclareMathOperator{\Poly}{\mathrm{Poly}}

\newcommand{\cL}{\mathcal{L}}
\newcommand{\cH}{\mathcal{H}}
\newcommand{\calP}{\mathcal{P}}
\newcommand{\PG}{\operatorname{PG}}
\newcommand{\Rep}{\operatorname{Rep}}
\newcommand{\gr}{\operatorname{gr}}
\newcommand{\VRep}{\operatorname{VRep}}
\newcommand{\Ind}{\operatorname{Ind}}
\newcommand{\Res}{\operatorname{Res}}
\newcommand{\Aut}{\operatorname{Aut}}
\newcommand{\col}{\operatorname{col}}
\newcommand{\asc}{\operatorname{asc}}
\newcommand{\bad}{\operatorname{bad}}
\newcommand{\des}{\operatorname{des}}
\newcommand{\luis}{\textcolor{NavyBlue}}
\newcommand{\qbinom}[2]{\begin{bmatrix} #1 \\ #2 \end{bmatrix}_q}

\newcommand{\calL}{\mathcal{L}}
\newcommand{\calH}{\mathcal{H}}
\newcommand{\bz}{\mathbb{Z}}
\newcommand{\lt}{\normalfont\text{lt}}
\newcommand{\calA}{\mathcal{A}}
\newcommand{\calS}{\mathcal{S}}
\newcommand{\ind}{\text{\normalfont Ind}}
\newcommand{\stab}{\text{\normalfont Stab}}
\newcommand{\chow}{\normalfont\text{\underline{H}}}
\newcommand{\augchow}{\normalfont\text{H}}
\newcommand{\klpoly}{\normalfont\text{P}}
\newcommand{\invklpoly}{\normalfont\hat{\text{Q}}}
\newcommand{\redchar}{\overline{\chi}}
\newcommand{\flats}{\normalfont\text{Flats}}
\newcommand{\flags}{\normalfont\text{Flags}}
\newcommand{\matM}{\normalfont\text{M}}
\newcommand{\hilbMmodF}{\normalfont\text{\underline{H}}_\text{M/F}}
\newcommand{\chibar}{\overline{\chi}}
\newcommand{\E}{\normalfont\text{E}}
\newcommand{\ZZ}{\mathbb{Z}}
\newcommand{\Ebasis}[1]{\mathrm{#1}}
\newcommand{\Cbasis}[1]{\mathrm{c}^{#1}}
\newcommand{\CbasisSub}[2]{\mathrm{c}^{#1}_{#2}}
\newcommand{\klbasis}[1]{\mathrm{z}^{#1}}
\newcommand{\klbasisSub}[2]{\mathrm{z}^{#1}_{#2}}
\newcommand{\Ibasis}[2]{\mathfrak{b}^{#1}_{#2}}
\newcommand{\pal}{\normalfont\text{Pal}}
\newcommand{\klchowcoeff}{\mathfrak{a}}
\newcommand{\IntP}{\normalfont\text{Int}}


%\newcommand\rev[1]{#1^{\operatorname{rev}}}

\AtBeginDocument{%
   \def\MR#1{}
}


%% preliminary version only
\begin{comment}
 \usepackage{prelim2e}
 \renewcommand{\PrelimText}{\sf\scriptsize --- INCOMPLETE DRAFT of \today ---}
 \marginparwidth=0.75in
 \usepackage[notcite,notref]{showkeys}
 \usepackage[colorinlistoftodos,bordercolor=orange,backgroundcolor=orange!20,linecolor=orange,textsize=scriptsize]{todonotes}

 \usepackage{prelim2e}
 \usepackage[us,24hr]{datetime}
 \renewcommand{\PrelimText}{\sf\scriptsize ---\,Version of \today \,-\,\currenttime\,---}
\end{comment}



\title[Deletion Formulas for Chow Functions]{Deletion Formulas for Chow Functions of ??}

\author[N.~Nepal]{Nutan Nepal}

\address{(N. Nepal)
Department of Mathematics, North Carolina State University, Raleigh, NC, USA.
}
\email{nnepal2@ncsu.edu}

\subjclass[2020]{Primary: 05B35}

\allowdisplaybreaks

\begin{document}

\begin{abstract}

\end{abstract}

\maketitle

\section{Introduction}

\subsection{Main results}

\section{Preliminaries}

\subsection{Posets and Incidence Algebras}
Let \(\calP\) be a \emph{locally finite, weakly ranked, partially ordered set}.
We denote by \(\IntP(\calP)\) the set of all closed intervals of \(\calP\) and
by \(\rho: \IntP(\calP) \to \ZZ\)
the \emph{weak rank function} satisfying \(\rho_{st} = \rho_{sw} + \rho_{wt}\) for any \(s \le w \le t\).

The \emph{incidence algebra} of \(\calP\) over \(\ZZ[x]\), denoted \(\mathscr{I}\), is the free \(\ZZ[x]\)-module on \(\IntP(\calP)\). An element \(a \in \mathscr{I}\) is a function that assigns a polynomial \(a_{st}(x)\) to each interval \([s,t]\). The algebra structure is given by the convolution product:
\[ (ab)_{su} := \sum_{s\leq t\leq u} a_{st}(x) \cdot b_{tu}(x), \qquad
    \text{ for every \(s\leq u\) in \(\calP\).}
\] The algebra is associative with a multiplicative identity \(\delta\).
An element \(a \in \mathscr{I}\) is invertible if and only if \(a_{ss}(x)\) is a unit
for every \(s \in \calP\). Two fundamental elements are the \emph{zeta function},
\(\zeta_{st}(x) = 1\), and its inverse, the \emph{Möbius function}, \(\mu = \zeta^{-1}\).

\subsection{Kernels and the KLS theory}
The weak rank function \(\rho\) induces a subalgebra \(\mathscr{I}_\rho \subseteq \mathscr{I}\) consisting of elements \(a\) such that \(\deg(a_{st}) \le \rho_{st}\) for every interval \([s,t]\). This subalgebra is equipped with an involution \(a \mapsto a^{\rev}\) defined by:
\[ (a^{\rev})_{st}(x) := x^{\rho_{st}} a_{st}(x^{-1}) \]
An element \(\kappa \in \mathscr{I}_\rho\) is a \emph{\((\calP,\rho)\)-kernel} if
\(\kappa_{ss}(x) = 1\) and it satisfies the property \(\kappa^{-1} = \kappa^{\rev}\).
The primary example is the \emph{characteristic function} \(\chi := \mu \cdot \zeta^{\rev}\),
which is a kernel for any weakly ranked poset.

As formulated in \cite{ferroni-matherne-vecchi}, a kernel uniquely determines two key functions:
\begin{theorem}\label{thm:kernel-properties} [Kazhdan-Lusztig-Stanley Functions]
    Let \(\kappa \in \mathscr{I}_\rho\) be a \((\calP,\rho)\)-kernel. There exist unique elements \(f, g \in \mathscr{I}_\rho\) (the right and left KLS functions, respectively) satisfying:
    \begin{enumerate}[\normalfont (i)]
        \item \(f_{ss}(x) = 1\) and \(g_{ss}(x) = 1\) for all \(s \in \calP\).
        \item \(\deg f_{st}(x) < \frac{1}{2} \rho_{st}\) and \(\deg g_{st}(x) < \frac{1}{2} \rho_{st}\) for all \(s < t\).
        \item \(f^{\rev} = \kappa \cdot f\) and \(g^{\rev} = g \cdot \kappa\).
    \end{enumerate}
\end{theorem}

\subsection{Chow Functions and Related Invariants}
From a kernel \(\kappa\) and its KLS functions, one can define several other
important polynomial invariants.
Central to this paper is the \emph{Chow function},
which is formulated in \cite{ferroni-matherne-vecchi} as follows:

\begin{definition}\label{def:chow-reduced} [Reduced Kernel and Chow Function]
    Let \(\kappa\) be a \((\calP,\rho)\)-kernel.
    \begin{enumerate}
        \item The \emph{reduced kernel} \(\overline{\kappa} \in \mathscr{I}_\rho\) is given by:
              \[ \overline{\kappa}_{st}(x) = \begin{cases} \frac{\kappa_{st}(x)}{x - 1} & \text{if } s < t, \\ -1 & \text{if } s = t.\end{cases}\]
        \item The \emph{\(\kappa\)-Chow function} \(\chow \in \mathscr{I}_\rho\) is defined as:
              \[ \chow = - (\overline{\kappa})^{-1}\]
    \end{enumerate}
\end{definition}

The theory also provides three related symmetric functions.

\begin{definition}\label{def:symmetric-invariants} [Symmetric Invariants]
    Let \(\kappa\) be a \((\calP,\rho)\)-kernel with KLS functions \(f\) and \(g\), and Chow function \(\chow\).
    \begin{enumerate}
        \item The \emph{Z-function} is \(Z := g^{\rev} \cdot f = g \cdot f^{\rev}\).
        \item The \emph{right augmented Chow function} is \(F := \chow \cdot f^{\rev}\).
        \item The \emph{left augmented Chow function} is \(G := g^{\rev} \cdot \chow\).
    \end{enumerate}
\end{definition}

These functions satisfy fundamental symmetry properties that we will use throughout this work. For any interval \([s,t]\) with \(s<t\):
\begin{itemize}
    \item \(\chow_{st}(x) = x^{\rho_{st} - 1} \chow_{st}(x^{-1})\).
    \item \(Z_{st}(x) = x^{\rho_{st}} Z_{st}(x^{-1})\).
    \item \(F_{st}(x) = x^{\rho_{st}} F_{st}(x^{-1})\) and \(G_{st}(x) = x^{\rho_{st}} G_{st}(x^{-1})\).
\end{itemize}

\section{The Module of Poset Elements and Perverse Subgroups}

Let \(R := \ZZ[x, x^{-1}]\) be the ring of Laurent polynomials with integer coefficients.
For any integer \(n\), we denote by \(\pal(n)\) the additive group of polynomials \(f(x) \in R\)
satisfying \(f(x) = x^n f(x^{-1})\). While \(\pal(0)\) is a subring of \(R\), we note that \(\pal(n)\)
for \(n \neq 0\) is not, as \(\pal(n) \cdot \pal(m) \subseteq \pal(n+m)\).

Let \(\calP\) be a locally finite, weakly ranked poset. We assume \(\calP\) to be bounded
with the minimum element \(\hat{0}\) and the maximum element \(\hat{1}\),
which allows for the definition of a global rank function \(\rho(t) := \rho_{\hat{0},t}\).
Let \(\kappa\) be a \((\calP, \rho)\)-kernel with left KLS function \(g\).

\begin{definition}\label{def:module-of-poset-elements}
    The module \(\calH(\calP)\) is the free \(R\)-module with a standard basis
    \(\{\Ebasis{s}\}_{s \in \calP}\). An element \(\alpha \in \calH(\calP)\) is a formal sum
    \[ \alpha = \sum_{s \in \calP} \alpha_s \cdot\Ebasis{s}, \quad \text{where } \alpha_s \in R. \]
\end{definition}

\subsection{The Perversity Condition via Convolution}
The notion of a perverse element in this work is a symmetry condition
motivated by the axiomatic properties of the
KLS and Chow functions in the Definition \ref{def:symmetric-invariants}.

\begin{definition}
    Let \(g\) be the left KLS function associated with a kernel \(\kappa\). For any \(s \in \calP\), we define the \(R\)-linear map \(S_s: \calH(\calP) \to R\) by
    \[ S_s(\alpha) := \sum_{t \ge s} g^{\rev}_{st}(x^{-1}) \cdot \alpha_t. \]
\end{definition}

This definition is crafted so that for the characteristic kernel (\(\kappa=\chi\)),
it recovers the operator \(S_s(\alpha) = \sum_{t \ge s} x^{\rho(s)-\rho(t)} \alpha_t\).

In \cite{braden-vysogorets}, the authors define the subgroup of \textbf{perverse elements}
\(\calH_p(\calP) \subseteq \calH(\calP)\). We adapt their definition to define a
larger subgroup \(\calH'_p(\calP) \subseteq \calH(\calP)\) as follows:
\begin{definition}
    The subgroup \(\calH'_p(\calP) \subseteq \calH(\calP)\) consists of all elements \(\alpha\) that satisfy
    \[ S_s(\alpha) \in \pal(0) \quad \text{for all } s \in \calP. \]
\end{definition}

\subsection{The Chow and Kazhdan-Lusztig Bases}

Let \(\kappa\) be a \((\calP, \rho)\)-kernel, with associated KLS functions \(f\) and \(g\),
and Chow function \(\chow\). Using these functions, we have the following
two important bases for the module \(\calH(\calP)\).

\begin{definition}\label{def:kazhdan-lusztig-basis}
    The \textbf{Kazhdan-Lusztig basis} \(\{\klbasis{t}\}_{t \in \calP}\) is defined by
    \[ \klbasis{t} = \sum_{s \le t} \klbasisSub{t}{s}\cdot\Ebasis{s} := \sum_{s \le t} g^{\rev}_{st}(x) \cdot f_{st}(x^{-2}) \cdot \Ebasis{s},
    \]
    where \(f\) is the right KLS function of \(\kappa\).
\end{definition}

\begin{definition}\label{def:chow-basis}
    The \textbf{Chow basis} \(\{\Cbasis{t}\}_{t \in \calP}\) is defined by
    \[
        \Cbasis{t} = \sum_{s \le t} \CbasisSub{t}{s}\cdot\Ebasis{s} :=
        \sum_{s \le t} g^{\rev}_{st}(x) \cdot \chow_{st}(x^{-2}) \cdot \Ebasis{s},
    \]
    where \(\chow\) is the \(\kappa\)-Chow function.
\end{definition}

\begin{remark}\label{rem:chow-basis-coefficients}
    From the properties of Chow functions, the coefficients have a special structure:
    \(\CbasisSub{t}{t} = 1 \in \pal(0)\), and \(\CbasisSub{t}{s} \in \pal(2)\) for all \(s < t\).
\end{remark}

\begin{definition}\label{def:palindromic-module}
    We define the module \(\calS(\calP)\) to be the free \(\pal(0)\)-module with basis \(\{\Cbasis{t}\}_{t \in \calP}\).
\end{definition}

\begin{proposition}\label{prop:chow-basis-in-Hp}
    For all \(t \in \calP\), \(\Cbasis{t} \in \calH'_p(\calP)\), i.e.,
    it satisfies \(S_s(\Cbasis{t}) \in \pal(0)\) for all \(s \in \calP\).
\end{proposition}
\begin{proof}
    For any \(s \in \calP\), if \(s \not\le t\), then \(S_s(\Cbasis{t}) = 0\). If \(s \le t\), we have
    \begin{align*}
        S_s(\Cbasis{t}) & = \sum_{v: s \le v \le t} g^{\rev}_{sv}(x^{-1})\cdot\CbasisSub{t}{v}                          \\
                        & = \sum_{v: s \le v \le t} g^{\rev}_{sv}(x^{-1})\cdot g^{\rev}_{vt}(x)\cdot \chow_{vt}(x^{-2}) \\
                        & =^? g^{\rev}_{st}(x)\cdot \sum_{v: s \le v \le t} g^{\rev}_{sv}(x^{-2})\cdot \chow_{vt}(x^{-2})\\
                        & = g^{\rev}_{st}(x)\left(g^{\rev}\chow\right)_{st}(x^{-2}) = g^{\rev}_{st}(x)G_{st}(x^{-2})
    \end{align*}
    We have that \(g^{\rev}_{st}(x)\in \pal(2\rho_{st})\) and \(G_{st}(x^{-2})\in \pal(-2\rho_{st})\). Their product is therefore in \(\pal(0)\), as required.
\end{proof}

\begin{proposition}\label{prop:Hp_is_pal0_mod}
    The subgroup \(\calH'_p(\calP)\) is a \(\pal(0)\)-module.
\end{proposition}
\begin{proof}
    Let \(\alpha = \sum_s p_s \cdot \Ebasis{s} \in \calH'_p(\calP)\) and \(q \in \pal(0)\).
    We see that
    \[ S_t(q \cdot \alpha) = \sum_{v \ge t} x^{\rho(t)-\rho(v)} (q \cdot p_v)
        = q \cdot S_t(\alpha) \in \pal(0), \]
    since \(\pal(0)\) is a ring. Therefore, \(q \cdot \alpha \in \calH'_p(\calP)\).
\end{proof}

\begin{proposition}\label{prop:LPDecomposition}
    Any Laurent polynomial \(f\) can be uniquely written as a sum \(f = f_0 + f_2\) where
    \(f_0 \in \pal(0)\) and \(f_2 \in \pal(2)\).The projections are given by
    \[ f_0(x) = \frac{x^2 f(x^{-1}) - f(x)}{x^2 - 1} \quad \text{and}
        \quad f_2(x) = \frac{f(x) - f(x^{-1})}{1 - x^{-2}}. \]

\end{proposition}

\begin{proposition}\label{prop:SisHp}
    \(\calS(\calP) = \calH'_p(\calP)\) as \(\pal(0)\)-modules.
\end{proposition}
\begin{proof}
    Proposition \ref{prop:chow-basis-in-Hp} and \ref{prop:Hp_is_pal0_mod}
    show that \(\calS(\calP) \subseteq \calH'_p(\calP)\). We now show the reverse inclusion.

    Let \(\alpha = \sum_s \alpha_s \cdot \Ebasis{s}\) be in \(\mathcal{H}'_p(\calP)\). Since \(\{\Cbasis{s}\}\) is a basis for \(\mathcal{H}(\calP)\), \(\alpha\) has a unique expansion \(\alpha = \sum_{s \in \calP} q_s \cdot \Cbasis{s}\) for \(q_s \in R\). We must prove that each \(q_s \in \pal(0)\) by induction on the corank of \(s\).

    The change-of-basis formula is \(\alpha_s = \sum_{t \ge s} q_t \cdot \CbasisSub{t}{s}\).

    \textbf{Base Case:} For \(s=\hat{1}\) (the maximal element), this simplifies to \(q_{\hat{1}} = \alpha_{\hat{1}}\). Since \(\alpha \in \mathcal{H}'_p(\calP)\), we have \(S_{\hat{1}}(\alpha) = \alpha_{\hat{1}} \in \pal(0)\), establishing the base case.

    \textbf{Inductive Step:} Assume for all \(t\) with \(\rho(t) > \rho(s)\), we have \(q_t \in \pal(0)\).
    We know \(S_s(\alpha) \in \pal(0)\). We also have \(S_s(\alpha) = \sum_{v \ge s} x^{\rho(s)-\rho(v)} \alpha_v\). Substituting the basis change:
    \[ S_s(\alpha) = \sum_{v \ge s} x^{\rho(s)-\rho(v)} \left( \sum_{t \ge v} q_t \cdot \CbasisSub{t}{v} \right)
        = \sum_{t \ge s} q_t \left( \sum_{v: s \le v \le t} x^{\rho(s)-\rho(v)} \CbasisSub{t}{v} \right)
        = \sum_{t \ge s} q_t \cdot S_s(\Cbasis{t}). \]
    Isolating the \(t=s\) term gives \(S_s(\alpha) = q_s \cdot S_s(\Cbasis{s}) + \sum_{t > s} q_t \cdot S_s(\Cbasis{t})\). Since \(S_s(\Cbasis{s})=1\), we rearrange to get:
    \[ q_s = S_s(\alpha) - \sum_{t > s} q_t \cdot S_s(\Cbasis{t}). \]
    By the inductive hypothesis, each \(q_t\) is in \(\pal(0)\). By Proposition \ref{prop:cinHp}, each \(S_s(\Cbasis{t})\) is in \(\pal(0)\). Since \(\pal(0)\) is a ring, the entire sum is in \(\pal(0)\). Thus, \(q_s\) must also be in \(\pal(0)\).
\end{proof}

\begin{proposition}[Change of Basis Formula]\label{prop:changeofbasis}
    If \(\alpha = \sum_s \alpha_s \cdot\Ebasis{s} \in \mathcal{S}(\calP)\),
    then its expansion in the \(\Cbasis{s}\) basis is
    \(\alpha = \sum_s \overline{\alpha_s} \cdot\Cbasis{s}\), where \(\overline{\alpha_s}\) is the unique \(\pal(0)\) component of \(\alpha_s\) from the decomposition in Proposition~\ref{prop:LPDecomposition}.
\end{proposition}
\begin{proof}
    Let \(\alpha \in \mathcal{S}(\calP)\). By definition, \(\alpha = \sum_s q_s \cdot \Cbasis{s}\) where each \(q_s \in \pal(0)\). The change-of-basis formula gives:
    \[ \alpha_s = q_s \cdot \CbasisSub{s}{s} + \sum_{t > s} q_t \cdot \CbasisSub{t}{s} = q_s + \sum_{t > s} q_t \cdot \CbasisSub{t}{s}. \]
    We proceed by induction on the corank of \(s\). The base case for \(s=\hat{1}\) is \(q_{\hat{1}} = \alpha_{\hat{1}}\), and since \(q_{\hat{1}} \in \pal(0)\), we have \(q_{\hat{1}} = \overline{\alpha_{\hat{1}}}\).

    For the inductive step, assume \(q_t = \overline{\alpha_t}\) for all \(t>s\). For any \(t>s\), \(q_t \in \pal(0)\) and \(\CbasisSub{t}{s} \in \pal(2)\), so their product is in \(\pal(2)\). The sum \(\sum_{t > s} q_t \cdot \CbasisSub{t}{s}\) is therefore in \(\pal(2)\).
    The equation shows \(\alpha_s\) is the sum of \(q_s\) (which is in \(\pal(0)\)) and a term in \(\pal(2)\). By the uniqueness of the decomposition in Proposition~\ref{prop:LPDecomposition}, \(q_s\) must be the unique \(\pal(0)\) component of \(\alpha_s\).
\end{proof}

\bibliographystyle{amsalpha}
\bibliography{references}

\end{document}